\subsection{Eindeutig rückwärts zuordnen}
\label{subsec:real_functions_inverse_principle}
Eine Rückwärtszuordnung ist nur dann eine Funktion, wenn kein Ausgangswert
zu zwei verschiedenen Eingaben gehört. Zum Beispiel liefern $-2$ und $2$
beim Quadrieren denselben Wert $4$: Auf ganz $\R$ lässt sich das
Quadrieren nicht eindeutig umkehren.

\par\noindent\begin{minipage}{\linewidth}
\begin{DefinitionBox}{Funktionseigenschaft -- Injektivität und Umkehrfunktion}{real_functions_inverse}
  Eine Funktion $f:D\to\R$ heißt \emph{injektiv} \EnglishTerm{injective},
  wenn verschiedene Eingaben verschiedene Ausgangswerte liefern:
  \[
    x_1\ne x_2\quad\Longrightarrow\quad f(x_1)\ne f(x_2).
  \]
  Dann existiert auf ihrer \textcolor{green!55!black}{Wertemenge} ${\color{green!55!black}W}=f(D)$ die
  \emph{Umkehrfunktion} \EnglishTerm{inverse function}
  $f^{-1}:W\to D$, bestimmt durch
  \[
    y=f(x)\quad\Longleftrightarrow\quad x=f^{-1}(y).
  \]
\end{DefinitionBox}
\begin{ExampleBox}[before skip=0.3\baselineskip]{Eine Geradengleichung umkehren}{real_functions_inverse_affine}
  Für $f(x)=2x+1$ soll zu einem Ausgangswert $y$ die Eingabe bestimmt werden.
  \begin{ClozeArea}[34]{6}
    Wir lösen die Gleichung nach der bisherigen Eingabe auf:
    \[
      \begin{array}{r@{\quad}l@{\qquad}l}
        &y=2x+1&\TransformOp{-1}\\[0.7em]
        \Longleftrightarrow&y-1=2x&\TransformOp{:2}\\[0.7em]
        \Longleftrightarrow&\dfrac{y-1}{2}=x.&
      \end{array}
    \]
    Somit ist $f^{-1}(y)=(y-1)/2$, oder mit dem üblichen Variablennamen
    $f^{-1}(x)=(x-1)/2$.
  \end{ClozeArea}
\end{ExampleBox}
\end{minipage}


\begin{KeyPointBox}{Hinweg und Rückweg heben sich auf}{real_functions_inverse_identities}
  Die Definitions- und \textcolor{green!55!black}{Wertemenge} tauschen ihre Rollen:
  \[
    f^{-1}(f(x))=x\quad(x\in D),\qquad
    f(f^{-1}(y))=y\quad(y\in W).
  \]
  Die Schreibweise $f^{-1}$ bezeichnet die Umkehrfunktion,
  \emph{nicht} den Kehrwert $1/f$.
\end{KeyPointBox}

\begin{TheoremBox}{Monotonie sichert die Umkehrbarkeit}{real_functions_monotone_inverse}
  Eine auf einem Intervall streng monotone Funktion ist injektiv.
  Ist sie zusätzlich stetig, ist ihre Umkehrfunktion auf der \textcolor{green!55!black}{Wertemenge}
  ebenfalls stetig und streng monoton in derselben Richtung.
\end{TheoremBox}
\begin{ExampleBox}{Den Definitionsbereich passend einschränken}{real_functions_square_inverse}
  Wir beschränken $f(x)=x^2$ auf $[0,\infty)$. Dort ist $f$ streng
  monoton steigend und hat die \textcolor{green!55!black}{Wertemenge} $[0,\infty)$.
  \[
    y=x^2,\ x\ge0
    \quad\Longleftrightarrow\quad
    x=\sqrt y.
  \]
  Die Umkehrfunktion ist daher $\Cloze{4cm}{f^{-1}(x)=\sqrt x}$.
  Eine Einschränkung auf $(-\infty,0]$ würde dagegen den Rückweg
  $y\mapsto-\sqrt y$ liefern.
  \begin{center}
    \CachedTikzFigure{content/resources/chapters/real_functions/tikz/functions_inverse_reflection.tex}
  \end{center}
\end{ExampleBox}
\begin{KeyPointBox}{Graphisch werden die Koordinaten vertauscht}{real_functions_inverse_graph}
  Aus $(x,y)$ auf dem Graphen von $f$ wird $(y,x)$ auf dem Graphen von
  $f^{-1}$. Bei gleicher Skalierung beider Achsen entspricht dies einer
  Spiegelung an der Geraden $y=x$.
\end{KeyPointBox}
